%% toc: P1 8 8.1,8.2,8.5,8.6
\chapter{Notation, Binding, and Comparison of Formal Presentations}
\chaptermark{Notation and Binding}
\label{sec:p01-w02-p01-024-notation}
\index{notation!scope}\index{binding}\index{precedence}

The symbol registry of \cref{sec:p01-symbol-registry} fixes, for every glyph,
a role, a scope and, for an operator, a rank.  This chapter supplies the
parts of that convention which a reader needs in order to use the registry
correctly and which are stated nowhere else on the reading route: the
registered uses of the Greek, Fraktur and blackboard alphabets by domain,
with the type-metavariable scopes; the occurrence-level definition of free
and bound variables; association and fixity for every operator that shares
one of the extended ranks; the difference between the atomic diagram and
the positive diagram; and the exact arity convention under which single
connectives are functionally complete, together with Post's criterion.

\emph{Inputs.}  The registry tables of \cref{sec:p01-symbol-registry}; the
first-order syntax of \cref{sec:fol-syntax-semantics-03}; the diagram
definitions of \cref{sec:diagrams}; the truth-function vocabulary of
\cref{sec:prop-syntax-semantics-04}; and the sixteen binary truth functions
of Table~\ref{tab:p01-sixteen}.

\emph{Conclusions.}  \cref{thm:p01-w02-p01-024-fv-occurrence} identifies
the recursive free-variable set with the set of variables having a free
occurrence.  \cref{thm:p01-w02-p01-024-positive-diagram} shows that the
positive diagram characterizes homomorphisms, while the full atomic diagram
characterizes embeddings.  \cref{thm:p01-w02-p01-024-nand-nor} proves that
NAND and NOR are each functionally complete for positive arities, and
\cref{thm:p01-w02-p01-024-no-closed} proves that no connective set without a
nullary member represents a nullary truth function.
\cref{thm:p01-w02-p01-024-post} is Post's criterion for positive-arity
truth functions, with the five maximal clones.  Every proof is in the
annex.

\section{Letter families and their registered uses by domain}
\label{sec:p01-w02-p01-024-letters}
\index{Greek letters}\index{Fraktur letters}\index{blackboard bold}

A typeface never assigns a role by itself; every role below comes from a
recorded symbol row, and the scope printed with it is the scope that row
records.  The complete extracts of \texttt{registry/symbols.csv} for the
snapshot identified in \cref{sec:p01-w02b-p01-012-extracts} are
Tables~\ref{tab:p01-w02b-p01-012-greek}, \ref{tab:p01-w02b-p01-012-fraktur},
\ref{tab:p01-w02b-p01-012-blackboard} and~\ref{tab:p01-w02b-p01-012-parts12}.
The lookup tables below point to those full records without repeating an
older extract.  The Greek-metavariable table provides a selected overview,
and the Parts~1--2 table retains the local declarations introduced with
this chapter.  Each scope remains the recorded part and chapter key, part
range, declaration label, range of labels, or \emph{global}; none of the
lookup routes identifies roles merely because their glyphs agree.

\begin{hbconvention}{def:p01-w02-p01-024-register-mentions}{Mentioning a scoped notation}
A letter displayed in a notation table, or in the comparison of letter
families below, is a specimen of the notation whose role and scope are
printed beside it.  The display does not extend that role to the surrounding
chapter.  Thus the blackboard-letter comparison can mention a truth-value
set, a Boolean algebra, a forcing order and a probability measure without
identifying their types.  The rank table likewise mentions an operator in
its declared source scope.  Actual use in a construction still requires
the local type declaration: in the functional-completeness discussion,
for example, the blackboard truth-value letter denotes the two-element
set of truth values.
\end{hbconvention}

\begin{hbconvention}{def:p01-w02-p01-024-language-filter-readings}{Two early local readings of letters}
In the logical-system triple of \cref{sec:levels-01}, \(\Lang\) is a
formal language: a set of finite expressions determined by its alphabet
and grammar.  That is the reading recalled by the opening logical-system
triple of \cref{sec:p01-w02-p01-003-language-change}.  After that section
explicitly fixes the first-order input, its language letter denotes the
signature of function and relation symbols with arities.  A signature
specifies the nonlogical vocabulary from which the corresponding terms
and formulas are generated; it is not itself the set of those expressions.
Similarly, \(\mathcal F\) in \cref{def:ultrafilter-lemma} is a proper
filter of subsets of the displayed underlying set: it contains the whole
set, excludes the empty set, and is closed under binary intersections and
supersets.  It is not the function-symbol family in
\cref{sec:p01-w02-p01-003-system-over-signature}.
These readings follow the explicit definitions in the two passages;
the common letter supplies no identification of their types.
\end{hbconvention}

\begin{hbconvention}{def:p01-w02-p01-024-inherited-domain-bars}{Domain bars in the inherited notation register and substitution problem}
In the inherited register \cref{app:notation}, the expression
\(\abs{\mathcal M}\) for a structure is explicitly its domain, a set.
The assignment row there consequently takes values in that set.
The same reading applies to the structure bars in the substitution
solution of \cref{sec:qualifying-problems-01}: the quantified elements
are values available to a variable assignment in the fixed structure.
In these two passages the bars return the carrier set, not its cardinal
number.  In the editorial notation this set is \(\dom{\mathcal M}\).
This convention records those two inherited readings and this explanatory
display; it leaves the cardinality notation for sets and the absolute-value
notation for numbers in their own declared scopes.
\end{hbconvention}

\begin{hbconvention}{def:p01-w02-p01-024-type-metavariables}{Greek ordinal, type and formula metavariables}
The lowercase letters \(\alpha,\beta,\gamma\) have no global role.  Each
reading below holds only inside the scope named with it:
\begin{enumerate}
\item type metavariables \(\alpha,\beta\) (with \(\gamma\) in the registry
  row) for the simple type theory declared at the head of
  \cref{sec:simple-type-theory-01}, together with the type former
  \(\alpha\to_{\mathrm{ty}}\beta\), the ascription \(A:\alpha\) and the
  typed abstraction \(\lambda x{:}\alpha.B\) recorded for the range
  \emph{Legacy pre-batch4 P2.3--P7.6}.  The registry row for the three
  letters records the scope \emph{P2 prop-metatheory}, whose reader region
  is the propositional metatheory chapter beginning at
  \cref{sec:prop-metatheory-01}; no type is declared in that chapter, and
  its sections use \(\alpha,\beta,\gamma\) as formula metavariables (item~6).
  The declaration of \cref{sec:simple-type-theory-01} is the one a reader of
  typed terms should apply;
\item type metavariables \(\alpha,\beta,\gamma\) in the scope
  \emph{P7 normalization-curry-howard to lambda-calculi-and-type-theories},
  and finite-type metavariables \(\alpha,\beta\) in the scope
  \emph{Legacy pre-batch4 P7.6 reader meaning};
\item type metavariables \(\alpha,\beta,\gamma\) in the scope
  \emph{P12 research-workflow}, with \(\alpha\to_{\mathrm{ty}}\beta\) and
  \(M{:}\alpha\) recorded for the same scope;
\item ordinal metavariables \(\alpha,\beta\) and, as additional ordinal
  metavariables, \(\gamma,\delta,\xi\) in the scope
  \emph{P6 zf-zfc to determinacy}; ordinal metavariables \(\alpha,\beta\)
  and \(\gamma\) in the scope \emph{Legacy pre-batch4 P7.5 reader
  meaning}; countable ordinals \(\alpha,\beta,\lambda\) in the scope
  \emph{P3 Morley countable model dependencies};
\item ordinal metavariables \(\alpha,\beta\) and the cardinal letter
  \(\kappa\) in Part~1, from \cref{def:ordinal-cardinal} through
  \cref{sec:set-prerequisites-06}; in Part~2 the same letters, with
  \(\xi,\lambda,\mu\), are bound to ordinals and cardinals inside single
  constructions, as in the enumeration \(\langle\varphi_\alpha:\alpha<\lambda\rangle\)
  and the witnesses \(d_\xi\) (\(\xi<\kappa\)) of
  \cref{sec:fol-metatheory-01} and the constants \(c_\alpha\)
  (\(\alpha<\kappa\)) of \cref{sec:fol-cardinality-normalforms-02};
\item propositional formula metavariables \(\alpha,\beta,\gamma\) in
  Part~2: in the substitution, functional-completeness and Tseitin
  passages of the propositional syntax chapter
  (\cref{sec:p02-propositional-substitution,sec:prop-syntax-semantics-04,sec:p02-tseitin-transformation}),
  in the translation \(\operatorname{tr}_{\mathsf H}\) of
  \cref{def:p02-hilbert-translation}, in the closure condition of
  \cref{def:p02-cpl-tableau}, in the NNF grammar of
  \cref{def:p02-nnf-dual}, over \(\operatorname{Form}(P)\) throughout
  \cref{sec:p02-free-boolean-completeness}, and over quantifier-free NNF
  matrices in \cref{def:p02-matrix-clause-construction};
\item formula metavariables \(\alpha\) (Davies--Humberstone, scope
  \emph{P8 modal-temporal-epistemic}) and path-formula metavariables
  \(\alpha,\beta\) (CTL-star, scope \emph{P8 conditionals-deontic-free});
  error targets \(\alpha,\beta\) in the scope \emph{P15 sequential
  evidence}.
\end{enumerate}
Two further uses of the letter are names, not metavariables: the
subscript of the alpha-equivalence relation \(\varphi=_{\alpha}\psi\) and
the rule label \(\alpha\) of the first-order deduction rules in
\cref{sec:fol-deduction-01}.  Apart from the registry scope discussed in
item~1, the passages named in different items do not overlap, so within any
one passage at most one reading is available and the passage selects it.  A
use of \(\alpha,\beta\) or \(\gamma\) outside all of these passages needs
its own scoped declaration before the letter is used.
\end{hbconvention}

The type-metavariable reading in Part~2 is therefore a declared role, not a
borrowing from Part~7: \cref{sec:simple-type-theory-01} declares
\(\alpha,\beta\) for simple types and states that they do not keep the
earlier ordinal role.  The earlier roles in Parts~1 and~2 are the ordinal
and cardinal letters of item~5 and the formula metavariables of item~6.

\subsection{Greek metavariable families}
\label{sec:p01-w02-p01-024-greek}

Table~\ref{tab:p01-w02-p01-024-greek} gives selected Greek metavariable,
context and substitution roles from an earlier registry extract, grouped
by the part that opens the recorded scope.  The complete snapshot is
Table~\ref{tab:p01-w02b-p01-012-greek}.  These are the rows in which a Greek letter ranges over
formulas, sentences, types, ordinals, contexts or substitution maps rather
than naming one fixed object.  The same letter recurs with unrelated roles:
\(\sigma\), for example, is a global sentence metavariable, a uniform
substitution in the scope \emph{P2 classical rejection construction}, a
substitution on free variables in the scope \emph{P7 classical cut
elimination} and an MLTT substitution in the scope \emph{P7
lambda-calculi-and-type-theories}.  Only the recorded scope decides which
reading a passage uses.  The last group of the table lists the
Greek formula, ordinal and cardinal metavariables that the text of
Parts~1 and~2 uses without a registry row of their own
(\cref{def:p01-w02-p01-024-type-metavariables}, items~5 and~6).  Greek letters with one fixed meaning in one scope
(for instance \(\varepsilon_0\), \(\omega\) or \(\mu\) for a measure) are
listed with all other Greek uses in
Table~\ref{tab:p01-w02-p01-024-greek-domains}.

\begingroup\small
\begin{longtable}{@{}p{0.22\textwidth}p{0.32\textwidth}p{0.34\textwidth}@{}}
\caption{Greek metavariable, context and substitution families.}
\label{tab:p01-w02-p01-024-greek}\\
\toprule
Notation & Recorded role & Recorded scope \\
\midrule
\endfirsthead
\caption[]{Greek metavariable, context and substitution families (continued).}\\
\toprule
Notation & Recorded role & Recorded scope \\
\midrule
\endhead
\midrule\multicolumn{3}{r}{\emph{continued}}\\\endfoot
\bottomrule\endlastfoot
\multicolumn{3}{@{}l}{\emph{global scope}}\\
\(\varphi,\allowbreak \psi,\allowbreak \chi\) & formula metavariables & global \\
\(\sigma,\allowbreak \tau\) & sentence metavariables & global \\
\(\Gamma,\allowbreak \Delta,\allowbreak \Sigma,\allowbreak \Pi\) & formula contexts & global \\
\(\varphi[t/x]\) & capture-avoiding substitution & global \\
\multicolumn{3}{@{}l}{\emph{scopes recorded under P2}}\\
\(\theta,\allowbreak \rho,\allowbreak \eta\) & simultaneous term substitutions & Parts 2--12 from prop-\allowbreak{}deduction to mechanization \\
\(\upsilon\) & uniform propositional substitution & P2 prop-\allowbreak{}syntax-\allowbreak{}semantics \\
\(\varphi[\upsilon]\) & uniform-substitution instance & P2 prop-\allowbreak{}syntax-\allowbreak{}semantics \\
\(\operatorname{supp}_{\mathrm{var}}(\theta)\) & variable support of a term substitution & Parts 2--12 from prop-\allowbreak{}deduction to mechanization \\
\(t\theta,\allowbreak \theta\eta\) & simultaneous substitution action and composition & Parts 2--12 from prop-\allowbreak{}deduction to mechanization \\
\(\theta\preceq_{\mathrm{inst}}\rho\) & substitution instance preorder & Parts 2--12 from prop-\allowbreak{}deduction to mechanization \\
\(\varphi[U^n/X^n]\) & capture-avoiding relation-variable substitution & P2 prop-\allowbreak{}metatheory \\
\(\alpha,\allowbreak \beta,\allowbreak \gamma\) & type metavariables & P2 prop-\allowbreak{}metatheory \\
\ensuremath{\sigma} & Simultaneous uniform substitution on propositional variables & P2 classical rejection construction \\
\ensuremath{\sigma_v} & Countervaluation-coded simultaneous formula substitution & P2 classical rejection construction \\
\multicolumn{3}{@{}l}{\emph{scopes recorded under P3}}\\
\(\sigma,\allowbreak \pi_i\) & Skolem variable substitution and variable-disjoint copy renamings & P3 Lyndon 1959 preservation constructive import \\
\(\chi\) & Boolean substitution on atomic formulas in source chain argument & P3 Lyndon 1959 preservation \\
\(\Sigma\) & set of substitution instances of the Skolem matrix & P3 Lyndon 1959 preservation chain argument \\
\multicolumn{3}{@{}l}{\emph{scopes recorded under P5}}\\
\(s_\theta^+,\allowbreak s_\theta^-,\allowbreak s_\varphi^+\) & fixed numeral-substitution and negated-root descriptions & P5 arithmetic-\allowbreak{}coding to reverse-\allowbreak{}mathematics \\
\multicolumn{3}{@{}l}{\emph{scopes recorded under P6}}\\
\(\alpha,\allowbreak \beta\) & ordinal metavariables & P6 zf-\allowbreak{}zfc to determinacy \\
\(\nu,\allowbreak \nu_f,\allowbreak \zeta\) & forcing-name metavariables & P6 cardinals-\allowbreak{}choice \\
\(\kappa,\allowbreak \lambda,\allowbreak \mu\) & ordinal or cardinal metavariables as explicitly quantified & P6 zf-\allowbreak{}zfc to determinacy \\
\(\gamma,\allowbreak \delta,\allowbreak \xi\) & additional ordinal metavariables & P6 zf-\allowbreak{}zfc to determinacy \\
\multicolumn{3}{@{}l}{\emph{scopes recorded under P7}}\\
\(\alpha,\allowbreak \beta,\allowbreak \gamma\) & type metavariables & P7 normalization-\allowbreak{}curry-\allowbreak{}howard to lambda-\allowbreak{}calculi-\allowbreak{}and-\allowbreak{}type-\allowbreak{}theories \\
\(d(\varphi),\allowbreak d(\Gamma)\) & defined-negation formula and context translation & P7 lk-\allowbreak{}lj \\
\(\bigwedge\Gamma\) & finite context conjunction & P7 lk-\allowbreak{}lj \\
\(\bigvee\Delta\) & finite context disjunction & P7 lk-\allowbreak{}lj \\
\(\Theta_{\mathrm{ty}}\) & term-variable typing context & Legacy \texttt{pre-\allowbreak{}batch4} P7.\allowbreak{}3-\allowbreak{}P7.\allowbreak{}6 reader meaning \\
\(\Delta_{\mathrm{ty}}\) & System-F type-variable context & P7 lambda-\allowbreak{}calculi-\allowbreak{}and-\allowbreak{}type-\allowbreak{}theories \\
\(M[\alpha],\allowbreak \alpha[\beta/a],\allowbreak M[\alpha/a]\) & System-F type application and capture-avoiding type substitution & P7 lambda-\allowbreak{}calculi-\allowbreak{}and-\allowbreak{}type-\allowbreak{}theories \\
\(\Gamma\;\mathsf{ctx}\) & MLTT context-formation judgment & P7 lambda-\allowbreak{}calculi-\allowbreak{}and-\allowbreak{}type-\allowbreak{}theories \\
\(\sigma,\allowbreak \tau:\Delta\to\Gamma\) & MLTT substitutions with common domain and codomain & P7 lambda-\allowbreak{}calculi-\allowbreak{}and-\allowbreak{}type-\allowbreak{}theories \\
\(\alpha,\allowbreak \beta\) & ordinal metavariables & Legacy \texttt{pre-\allowbreak{}batch4} P7.\allowbreak{}5 reader meaning \\
\(\gamma\) & ordinal metavariable & Legacy \texttt{pre-\allowbreak{}batch4} P7.\allowbreak{}5 reader meaning \\
\(\alpha,\allowbreak \beta\) & finite-type metavariables & Legacy \texttt{pre-\allowbreak{}batch4} P7.\allowbreak{}6 reader meaning \\
\(\pi\) & encoded propositional proof metavariable & Legacy \texttt{pre-\allowbreak{}batch4} P7.\allowbreak{}8 reader meaning \\
\(\Theta,\allowbreak \Delta\) & finite source and target typing contexts & sec:\allowbreak{}p07-\allowbreak{}stlc-\allowbreak{}metatheory \\
\(\rho\) & simultaneous substitution of raw terms for free variables & sec:\allowbreak{}p07-\allowbreak{}stlc-\allowbreak{}metatheory \\
\(\Delta,\allowbreak \Gamma\) & type-variable and term-variable typing contexts & sec:\allowbreak{}p07-\allowbreak{}system-\allowbreak{}f-\allowbreak{}normalization \\
\(\tau\) & simultaneous capture-avoiding type substitution & sec:\allowbreak{}p07-\allowbreak{}system-\allowbreak{}f-\allowbreak{}normalization \\
\(\sigma\) & simultaneous capture-avoiding term substitution & sec:\allowbreak{}p07-\allowbreak{}system-\allowbreak{}f-\allowbreak{}normalization \\
\(\mathsf{WF}(\Delta,\allowbreak \alpha)\) & finite type-formation check in a type context & sec:\allowbreak{}p07-\allowbreak{}system-\allowbreak{}f-\allowbreak{}confluence-\allowbreak{}and-\allowbreak{}checking \\
\(\theta\) & finite simultaneous substitution on free variables & P7 intuitionistic substitution \\
\(\Gamma,\allowbreak \Delta,\allowbreak \Pi,\allowbreak \Sigma\) & retained finite multiset contexts & P7 classical cut elimination \\
\(\sigma\) & finite simultaneous substitution on free variables & P7 classical cut elimination \\
\(C\sigma\) & formula after a finite simultaneous substitution & P7 classical cut elimination \\
\(\sigma,\allowbreak \tau\) & simultaneous source formula substitutions & P7 Cook formula substitutions \\
\(\Gamma;\allowbreak \Delta,\allowbreak \Gamma[x:A]\) & source context concatenation and fresh final declaration & \texttt{sec:\allowbreak{}p07-\allowbreak{}ch88} \\
\(\Gamma_?\) & quest-prefixed auxiliary occurrence context & P7 girard-\allowbreak{}1987-\allowbreak{}cut-\allowbreak{}trips-\allowbreak{}and-\allowbreak{}boxes \\
\multicolumn{3}{@{}l}{\emph{scopes recorded under P8}}\\
\(\Gamma_0\) & finite-support witness context & P8 intuitionistic-\allowbreak{}logic \\
\(\alpha\) & Davies--Humberstone formula metavariable & P8 modal-\allowbreak{}temporal-\allowbreak{}epistemic \\
\(\Phi\) & Fitting quantified-modal formula metavariable & P8 modal-\allowbreak{}temporal-\allowbreak{}epistemic \\
\(\alpha,\allowbreak \beta\) & CTL-star path-formula metavariables & P8 conditionals-\allowbreak{}deontic-\allowbreak{}free \\
\(\Phi\) & ordered Lambek context metavariable & P8 paraconsistent-\allowbreak{}and-\allowbreak{}nonmonotonic \\
\multicolumn{3}{@{}l}{\emph{scopes recorded under P9}}\\
\(\operatorname{psub}_{P_i}(\phi,\allowbreak \psi),\allowbreak [\phi]^\psi_{P_i}\) & Kamp proper substitution for a unary predicate & P9 legacy unit 4 temporal-\allowbreak{}automata \\
\(\eta[\zeta/q_i]\) & single-pass temporal proposition substitution & P9 legacy unit 4 temporal-\allowbreak{}automata \\
\ensuremath{C^\delta} & Total state function induced by a ground context & p09 finite-\allowbreak{}tree-\allowbreak{}algebra annex \\
\multicolumn{3}{@{}l}{\emph{scopes recorded under P10}}\\
\(t\sigma,\allowbreak \widehat\sigma\) & simultaneous algebraic term substitution & P10 algebraic-\allowbreak{}and-\allowbreak{}categorical-\allowbreak{}logic \\
\(f_i(\chi)\) & one-coordinate modal formula context & P10 legacy unit 3 nonclassical-\allowbreak{}algebras \\
\(\mathcal S_{\beta\eta}\) & category of arrow-only typed contexts & P10 legacy unit 8 cartesian-\allowbreak{}lambda \\
\(\sigma:\Theta\to\Psi,\allowbreak M[\sigma]\) & typed simultaneous substitution and its action & P10 legacy unit 8 cartesian-\allowbreak{}lambda \\
\(c_\Gamma(\beta),\allowbreak \Delta^\Gamma\) & curried type and context exponential & P10 legacy unit 8 cartesian-\allowbreak{}lambda \\
\(\llbracket\Gamma\rrbracket,\allowbreak \llbracket M\rrbracket\) & categorical interpretation of a context or term & P10 legacy unit 8 cartesian-\allowbreak{}lambda \\
\(\underline A,\allowbreak A_T,\allowbreak P_\Gamma\) & named type interpreted type and context object & P10 legacy unit 8 cartesian-\allowbreak{}lambda \\
\(\operatorname{pack}_\Gamma,\allowbreak \operatorname{unpack}_\Gamma\) & inverse context packing arrows & P10 legacy unit 8 cartesian-\allowbreak{}lambda \\
\(\pi_{I,A},\allowbreak f\times1_A\) & context projection and parameter substitution & P10 legacy unit 9 doctrine-\allowbreak{}quantifiers \\
\(\Gamma,\allowbreak \Phi\) & variable context and hypothesis list in Pitts's calculus & P10 legacy unit 9 doctrine-\allowbreak{}quantifiers \\
\(\Phi\vdash\psi[\Gamma]\) & contextual predicate sequent & P10 legacy unit 9 doctrine-\allowbreak{}quantifiers \\
\(I_\Gamma\) & interpreted finite variable context & P10 legacy unit 9 doctrine-\allowbreak{}quantifiers \\
\(\llbracket M\rrbracket_\Gamma,\allowbreak \llbracket\varphi\rrbracket_\Gamma\) & contextual term and predicate interpretations & P10 legacy unit 9 doctrine-\allowbreak{}quantifiers \\
\(s:I_\Delta\to I_\Gamma\) & interpreted simultaneous term substitution & P10 legacy unit 9 doctrine-\allowbreak{}quantifiers \\
\([\phi]\) & extension subobject of a formula in its declared context & p10/\allowbreak{}n-\allowbreak{}d-\allowbreak{}023-\allowbreak{}source-\allowbreak{}03 \\
\multicolumn{3}{@{}l}{\emph{scopes recorded under P11}}\\
\(\operatorname{diag}_{\Pi}\) & effective self-index substitution on analytical formula codes & P11 Kripke truth fixed points \\
\multicolumn{3}{@{}l}{\emph{scopes recorded under P12}}\\
\(\alpha,\allowbreak \beta,\allowbreak \gamma\) & type metavariables & P12 research-\allowbreak{}workflow \\
\(\Theta_{\mathrm{ty}}\) & term-variable typing context & P12 research-\allowbreak{}workflow \\
\(\sigma\) & finite-support term substitution & sec:\allowbreak{}p12-\allowbreak{}small-\allowbreak{}clause-\allowbreak{}standard \\
\(\nu_\sigma\) & valuation induced by evaluating a substitution & annex:\allowbreak{}p12.\allowbreak{}small.\allowbreak{}clause.\allowbreak{}simplify \\
\(\sigma\) & capture-free substitution & sec:\allowbreak{}p12-\allowbreak{}opt-\allowbreak{}small-\allowbreak{}clauses \\
\(\theta,\allowbreak \delta\) & most general unifier and residual substitution in critical-overlap factorization & sec:\allowbreak{}p12-\allowbreak{}rewriting \\
\(\sigma,\allowbreak \sigma'\) & substitutions before and after changing one variable image & sec:\allowbreak{}p12-\allowbreak{}rewriting \\
\(\rho\) & combined matching substitution for renamed rule variables & annex:\allowbreak{}p12.\allowbreak{}rewriting.\allowbreak{}critical-\allowbreak{}pairs \\
\(\theta\eta\) & simultaneous substitution composition with theta acting first & sec:\allowbreak{}p12-\allowbreak{}unification-\allowbreak{}extraction-\allowbreak{}adequacy \\
\(\mu\) & finite-support most general substitution and the subsequently extracted answer & sec:\allowbreak{}p12-\allowbreak{}unification-\allowbreak{}extraction-\allowbreak{}adequacy \\
\(\mathrm{Type}_\Gamma(M),\allowbreak \mathrm{Ok}(\Gamma)\) & source synthesis and context-validation functions & \texttt{p12/\allowbreak{}har2001-\allowbreak{}c18} \\
\(t^\sigma\) & homomorphic substitution action & p12-\allowbreak{}src-\allowbreak{}027-\allowbreak{}501 \\
\(q\theta\rho\) & updated source skeleton with first-then-second finite substitution composition & P12 equational hereditary fragments \\
\(\sigma_\Pi\) & identification substitution & p12-\allowbreak{}equational-\allowbreak{}combination \\
\(\mathsf{sub}_\sigma,\allowbreak \nu_\sigma\) & syntax substitution and composed valuation & p12/\allowbreak{}har-\allowbreak{}autarkic-\allowbreak{}interpretation \\
\multicolumn{3}{@{}l}{\emph{scopes recorded under P13}}\\
\ensuremath{\Gamma} & Finite categorical premise context in editorial bookkeeping & p13 \texttt{src063-\allowbreak{}0103-\allowbreak{}deduction} \\
\multicolumn{3}{@{}l}{\emph{scopes recorded under P14}}\\
\(\Gamma^N,\allowbreak \Gamma^{\mathrm{Ko}},\allowbreak \Gamma^{\mathrm{Ku}}\) & memberwise translated finite formula contexts & P14 proof translation comparisons \\
\(\Gamma,\allowbreak \Delta\) & formula-premise sets or exact sheet-indexed resource contexts & P14 nonclassical modal and resource sheet comparisons \\
\(\Gamma,\allowbreak \Phi\) & distinct sorted-variable context and finite formula-hypothesis list & P14 nonclassical categorical sheet comparisons \\
\multicolumn{3}{@{}l}{\emph{scopes recorded under P15}}\\
\(I,\allowbreak J,\allowbreak T,\allowbreak \lambda,\allowbreak \chi_R,\allowbreak Q_F,\allowbreak x_{m,j},\allowbreak \mu_m,\allowbreak F_m,\allowbreak g,\allowbreak h\) & intervals monotone substitution Lebesgue measure continuous cutoffs quantile midpoint atoms approximating laws and controlling functions & P15 characteristic-\allowbreak{}function-\allowbreak{}criteria-\allowbreak{}and-\allowbreak{}moments \\
\multicolumn{3}{@{}l}{\emph{scopes recorded under P16}}\\
\(\phi,\allowbreak \psi,\allowbreak F\) & formula metavariables & sec:\allowbreak{}p16-\allowbreak{}nonmonotonic-\allowbreak{}causal-\allowbreak{}theories \\
\multicolumn{3}{@{}l}{\emph{scopes recorded under P17}}\\
\(\theta\) & consistent substitution used in grounding and method unification & P17.\allowbreak{}HTN \\
\multicolumn{3}{@{}l}{\emph{scopes recorded under P19}}\\
\(B,\allowbreak T,\allowbreak m,\allowbreak \tau_m\) & finite base and target object sets correspondence and induced expression substitution & P19 analogy-\allowbreak{}and-\allowbreak{}problem-\allowbreak{}transfer \\
\multicolumn{3}{@{}l}{\emph{conventions of Parts~1--2 recorded with this chapter}}\\
\(\alpha,\allowbreak \beta;\allowbreak \kappa\) & ordinal and cardinal metavariables & sec:\allowbreak{}set-\allowbreak{}prerequisites-\allowbreak{}05 through sec:\allowbreak{}set-\allowbreak{}prerequisites-\allowbreak{}06 \\
\(\alpha,\allowbreak \beta,\allowbreak \gamma\) & propositional formula metavariables & P2 prop-\allowbreak{}syntax-\allowbreak{}semantics \\
\(\alpha,\allowbreak \beta,\allowbreak \gamma\) & propositional formula metavariables & sec:\allowbreak{}p02-\allowbreak{}hilbert-\allowbreak{}cpl \\
\(\alpha\) & propositional formula metavariable & sec:\allowbreak{}p02-\allowbreak{}analytic-\allowbreak{}tableaux \\
\(\alpha,\allowbreak \beta\) & NNF formula metavariables & sec:\allowbreak{}p02-\allowbreak{}propositional-\allowbreak{}duality \\
\(\alpha,\allowbreak \beta,\allowbreak \gamma\) & formula metavariables over \(\operatorname{Form}(P)\) & P2 free Boolean syntax \\
\(\alpha,\allowbreak \beta\) & quantifier-free NNF matrix metavariables & sec:\allowbreak{}p02-\allowbreak{}fol-\allowbreak{}clausal-\allowbreak{}pipeline \\
\end{longtable}
\endgroup

\subsection{Every registered Greek use by domain}
\label{sec:p01-w02-p01-024-greek-domains}

Table~\ref{tab:p01-w02-p01-024-greek-domains} directs the reader to every
Greek-family row selected from the declared registry snapshot, including
the separate formula-metavariable rows for \(\alpha,\beta,\gamma\) in
Part~2.  The complete extract preserves each row's role and scope.

\begingroup\small
\begin{longtable}{@{}p{0.32\textwidth}p{0.60\textwidth}@{}}
\caption{Greek letters by domain: complete lookup route.}
\label{tab:p01-w02-p01-024-greek-domains}\\
\toprule
Lookup & Complete record \\
\midrule
Registered roles and scopes & Table~\ref{tab:p01-w02b-p01-012-greek}, using the snapshot and selection rule of \cref{def:p01-w02b-p01-012-selection}. \\
Local metavariable declarations & \cref{def:p01-w02-p01-024-type-metavariables} and Table~\ref{tab:p01-w02-p01-024-parts12}. \\
\bottomrule
\end{longtable}
\endgroup

\subsection{Fraktur letters}
\label{sec:p01-w02-p01-024-fraktur}

No Fraktur letter has a global row.  Table~\ref{tab:p01-w02-p01-024-fraktur}
directs the reader to the complete Fraktur extract, including Parts~1
and~2, for the snapshot identified in
\cref{sec:p01-w02b-p01-012-extracts}.  Each row confines its Fraktur letter to
the recorded scope; several unrelated roles (structures, frames, logic
metavariables, cardinal characteristics) are separated only by those
scopes.

\begingroup\small
\begin{longtable}{@{}p{0.32\textwidth}p{0.60\textwidth}@{}}
\caption{Fraktur letters by domain: complete lookup route.}
\label{tab:p01-w02-p01-024-fraktur}\\
\toprule
Lookup & Complete record \\
\midrule
Registered roles and scopes & Table~\ref{tab:p01-w02b-p01-012-fraktur}, using the snapshot and selection rule of \cref{def:p01-w02b-p01-012-selection}. \\
\bottomrule
\end{longtable}
\endgroup

\subsection{Blackboard letters}
\label{sec:p01-w02-p01-024-blackboard}

The global blackboard letters are \(\mathbb N\) (with \(0\in\mathbb N\)),
\(\mathbb Z\), \(\mathbb Q\) and \(\mathbb R\).  Every other blackboard
letter is scoped: \(\mathbb B\) is the CPL truth-value set under P2, a
complete Boolean algebra under P6, and \(\mathbb B_n(T)\) a Boolean
quotient of formulas under P3; \(\mathbb F_2\) is the two-element field
under P2 and P15; \(\mathbb P\) is a forcing poset under P6 and a
probability measure under P11, P15, P16, P17 and P19; \(\mathbb E\) is an
expectation under P15, P16 and P17; \(\mathbb F\) is the category of finite
ordinals under P10 and a reduction strategy under P12.  The scope groups
below are headed by the part identifier with which the recorded scope
begins.
Table~\ref{tab:p01-w02-p01-024-blackboard} directs the reader to the
complete blackboard extract for the snapshot identified in
\cref{sec:p01-w02b-p01-012-extracts}, including rows in which the letter
is a subscript or an argument.

\begingroup\small
\begin{longtable}{@{}p{0.32\textwidth}p{0.60\textwidth}@{}}
\caption{Blackboard letters by domain: complete lookup route.}
\label{tab:p01-w02-p01-024-blackboard}\\
\toprule
Lookup & Complete record \\
\midrule
Registered roles and scopes & Table~\ref{tab:p01-w02b-p01-012-blackboard}, using the snapshot and selection rule of \cref{def:p01-w02b-p01-012-selection}. \\
\bottomrule
\end{longtable}
\endgroup

\subsection{Connectives, quantifiers, turnstiles and metavariables of
Parts~1 and~2}
\label{sec:p01-w02-p01-024-parts12}

Table~\ref{tab:p01-w02-p01-024-parts12} retains the local conventions of
Parts~1 and~2 recorded with this chapter and gives the complete lookup route.
Every row of global scope, with a scope opening in Part~1 or Part~2, or selected
by the label exception of \cref{def:p01-w02b-p01-012-selection}, in the
snapshot identified there, is in
Table~\ref{tab:p01-w02b-p01-012-parts12}; this includes the term metavariables
\(r,t,u\), the formula metavariables \(\alpha,\beta,\gamma\) and \(A\),
which that snapshot records as family-member rows with the scopes of
the passages that use them.

A \emph{convention} in this table is a standing reading of a letter or
sign: one declared for a chapter or section, or one used schematically in a
definition, rule or table without a declaration.  A letter that a single
definition, theorem or proof binds where it is used (``let \(\kappa\) be an
infinite cardinal'', ``for every \(\gamma\in\operatorname{supp}(\Gamma)\)'')
carries its declaration with it and is not repeated here; item~5 of
\cref{def:p01-w02-p01-024-type-metavariables} lists the ordinal and
cardinal letters bound in this way in Part~2.  The alternative glyphs in the
last column of Table~\ref{tab:p01-sixteen} are literature variants, not
operators of the handbook, and are not rows.  The fourth column is the
precedence rank or descriptor; \emph{meta} marks a metalinguistic relation
that never occurs inside an object-language formula, and \emph{n/a} a
symbol that is not an operator.

A row with a narrower scope never overrides the global row outside that
scope.  For instance, \(\vdash_P\) and \(\models_P\) are the derivability and
consequence relations of the free Boolean construction, and
\(\vdash_{\mathsf S}\) is direct derivability in the syllogistic fragment,
each in its recorded scope only.  Likewise the bare \(\vdash\) of the scope
\emph{P2 classical rejection construction} is a closed assertion judgment
paired with the rejection sign \(\dashv\) of the same scope; outside that
scope the global row \(\Gamma\vdash_{S}\varphi\) applies, whose system
subscript may be dropped only under a declared local abbreviation, as on the
Hilbert calculus sheet used in \cref{def:p02-hilbert-derivation}.  The negated turnstile
\(\Gamma\nvdash_{S}\varphi\) says that no derivation of \(\varphi\) from
\(\Gamma\) exists in \(S\); it is the negation of the global row, not a
further consequence relation, and it is used for consistency
(\(\Gamma\nvdash_{S}\bottom\)) and in the contrapositive completeness
arguments of \cref{sec:prop-metatheory-03,sec:fol-metatheory-03}.

\begingroup\small
\begin{longtable}{@{}p{0.22\textwidth}p{0.30\textwidth}p{0.24\textwidth}p{0.12\textwidth}@{}}
\caption{Local notation declarations of Parts~1--2 and the complete lookup route.}
\label{tab:p01-w02-p01-024-parts12}\\
\toprule
Notation & Recorded role & Recorded scope & Prec. \\
\midrule
\endfirsthead
\caption[]{Local notation declarations of Parts~1--2 (continued).}\\
\toprule
Notation & Recorded role & Recorded scope & Prec. \\
\midrule
\endhead
\midrule\multicolumn{4}{r}{\emph{continued}}\\\endfoot
\bottomrule\endlastfoot
\multicolumn{4}{@{}p{0.92\textwidth}@{}}{The complete registered rows and their scopes are in Table~\ref{tab:p01-w02b-p01-012-parts12}; the local declarations retained below keep their stated scopes.}\\
\midrule
\(\alpha,\allowbreak \beta;\allowbreak \kappa\) & ordinal and cardinal metavariables & sec:\allowbreak{}set-\allowbreak{}prerequisites-\allowbreak{}05 through sec:\allowbreak{}set-\allowbreak{}prerequisites-\allowbreak{}06 & n/a \\
\(\operatorname{NAND},\allowbreak \operatorname{NOR}\) & binary connective symbols interpreted by rows 1110 and 1000 & sec:\allowbreak{}p01-\allowbreak{}w02-\allowbreak{}p01-\allowbreak{}024-\allowbreak{}post & n/a \\
\(\alpha,\allowbreak \beta,\allowbreak \gamma\) & propositional formula metavariables & P2 prop-\allowbreak{}syntax-\allowbreak{}semantics & n/a \\
\(\alpha,\allowbreak \beta,\allowbreak \gamma\) & propositional formula metavariables & sec:\allowbreak{}p02-\allowbreak{}hilbert-\allowbreak{}cpl & n/a \\
\(A,\allowbreak B\) & formula metavariables of the negation-implication language & P2 classical rejection construction & n/a \\
\(\alpha\) & propositional formula metavariable & sec:\allowbreak{}p02-\allowbreak{}analytic-\allowbreak{}tableaux & n/a \\
\(\Gamma\nvdash_{S}\varphi\) & negated derivability: no derivation in the named calculus & P2 prop-\allowbreak{}metatheory & meta \\
\(\alpha,\allowbreak \beta\) & NNF formula metavariables & sec:\allowbreak{}p02-\allowbreak{}propositional-\allowbreak{}duality & n/a \\
\(\alpha,\allowbreak \beta,\allowbreak \gamma\) & formula metavariables over \(\operatorname{Form}(P)\) & P2 free Boolean syntax & n/a \\
\(r,\allowbreak t,\allowbreak u\) & term metavariables & P2 fol-\allowbreak{}syntax-\allowbreak{}semantics & n/a \\
\(\Gamma\nvdash_{S}\varphi\) & negated derivability: no derivation in the named calculus & sec:\allowbreak{}fol-\allowbreak{}metatheory-\allowbreak{}01 through sec:\allowbreak{}fol-\allowbreak{}metatheory-\allowbreak{}03 & meta \\
\(A,\allowbreak B\) & Hilbert-fragment formula metavariables & sec:\allowbreak{}p02-\allowbreak{}hilbert-\allowbreak{}fol-\allowbreak{}adequacy & n/a \\
\(\alpha,\allowbreak \beta\) & quantifier-free NNF matrix metavariables & sec:\allowbreak{}p02-\allowbreak{}fol-\allowbreak{}clausal-\allowbreak{}pipeline & n/a \\
\end{longtable}
\endgroup

\section{Free and bound occurrences}
\label{sec:p01-w02-p01-024-binding}
\index{variable!free occurrence}\index{variable!bound occurrence}

\Cref{sec:fol-syntax-semantics-03} defines the set \(\FV(\varphi)\) by
recursion and describes a bound occurrence informally.  The definitions
below make each occurrence-level notion exact, so that statements such as
``\(x\) occurs both free and bound in \(\varphi\)'' have a definite meaning.
The language is the first-order language of that section: terms are
variables and \(f(t_1,\ldots,t_n)\) for an \(n\)-ary function symbol
(\(n=0\) for a constant); formulas are \(r=t\), \(R(t_1,\ldots,t_n)\),
\(\bottom\), \(\neg\varphi\), \(\varphi\circ\psi\) for each binary
connective \(\circ\) of the language, \(\forall x\,\varphi\) and
\(\exists x\,\varphi\).  By unique readability each formula has exactly one
of these forms, so the recursions below are well defined.

\begin{hbdefinition}{def:p01-w02-p01-024-addresses}{Addresses and variable occurrences}
An \emph{address} is a finite sequence of natural numbers; \(\epsilon\) is
the empty address, and \(i\cdot a\) is the address \(a\) with \(i\) placed in
front.  An address \(b\) is a \emph{prefix} of \(a\), written
\(b\sqsubseteq a\), when \(a=b\cdot c\) for some address \(c\).  For a term
or formula \(e\), the set \(\operatorname{VO}(e)\) of
\emph{variable-occurrence addresses}, and the variable
\(e[a]\) at each \(a\in\operatorname{VO}(e)\), are defined by recursion:
\begin{enumerate}
\item \(\operatorname{VO}(x)=\{\epsilon\}\) and \(x[\epsilon]=x\);
\item \(\operatorname{VO}(f(t_1,\ldots,t_n))=
  \{i\cdot a:1\leq i\leq n,\ a\in\operatorname{VO}(t_i)\}\), and the same
  clause for \(R(t_1,\ldots,t_n)\) and for \(r=t\) (with \(t_1=r\),
  \(t_2=t\)), where \((\cdots)[i\cdot a]=t_i[a]\);
\item \(\operatorname{VO}(\bottom)=\varnothing\);
  \(\operatorname{VO}(\neg\varphi)=\{1\cdot a:a\in\operatorname{VO}(\varphi)\}\);
  \(\operatorname{VO}(\varphi\circ\psi)=
  \{1\cdot a:a\in\operatorname{VO}(\varphi)\}\cup
  \{2\cdot a:a\in\operatorname{VO}(\psi)\}\), with the variable read off in
  the immediate subformula;
\item for \(Q\in\{\forall,\exists\}\),
  \(\operatorname{VO}(Qx\,\varphi)=\{0\}\cup
  \{1\cdot a:a\in\operatorname{VO}(\varphi)\}\), where
  \((Qx\,\varphi)[0]=x\) and \((Qx\,\varphi)[1\cdot a]=\varphi[a]\).
\end{enumerate}
An \emph{occurrence of} \(x\) in \(e\) is an address \(a\) with
\(e[a]=x\).  The address \(0\) of a quantified formula is the
\emph{binding occurrence} written after the quantifier.  The set
\(\operatorname{QO}(\varphi)\) of \emph{quantifier-occurrence addresses} is
\(\varnothing\) for atomic formulas and \(\bottom\); \(\{1\cdot b:b\in
\operatorname{QO}(\varphi)\}\) for \(\neg\varphi\); the union of the
prefixed sets of both components for \(\varphi\circ\psi\); and
\(\{\epsilon\}\cup\{1\cdot b:b\in\operatorname{QO}(\varphi)\}\) for
\(Qx\,\varphi\).  The \emph{subformula occurrence} \(\varphi\{b\}\) at
\(b\in\operatorname{QO}(\varphi)\) is given by \(\varphi\{\epsilon\}=\varphi\),
\((\neg\psi)\{1\cdot b\}=\psi\{b\}\),
\((\psi_1\circ\psi_2)\{i\cdot b\}=\psi_i\{b\}\) and
\((Qx\,\psi)\{1\cdot b\}=\psi\{b\}\).  Each \(\varphi\{b\}\) has the form
\(Qy\,\chi\), and the quantifier occurrence at \(b\) \emph{quantifies} the
variable \(y\).
\end{hbdefinition}

The address \(0\) makes the variable written after a quantifier an
occurrence even when the quantifier is vacuous, as in \(\forall x\,P(y)\),
in agreement with the set \(\operatorname{Var}(\varphi)\) of
\cref{sec:fol-syntax-semantics-03}, whose quantifier clause adds the bound
variable.

\begin{hbdefinition}{def:p01-w02-p01-024-free-bound}{Scope, binding, free and bound occurrences}
Let \(\varphi\) be a formula, \(b\in\operatorname{QO}(\varphi)\) and
\(a\in\operatorname{VO}(\varphi)\).
\begin{enumerate}
\item The \emph{scope} of the quantifier occurrence \(b\) is the set of
  variable-occurrence addresses \(a\) with \(b\sqsubseteq a\).
\item The quantifier occurrence \(b\) \emph{binds} the occurrence \(a\) iff
  \(a\) is in the scope of \(b\), \(b\) quantifies the variable
  \(\varphi[a]\), and no quantifier occurrence \(b'\neq b\) with
  \(b\sqsubseteq b'\sqsubseteq a\) quantifies \(\varphi[a]\).
\item The occurrence \(a\) is \emph{bound} in \(\varphi\) iff some
  quantifier occurrence binds it, and \emph{free} in \(\varphi\)
  otherwise.
\item A variable \(x\) is \emph{free in} \(\varphi\) iff some occurrence of
  \(x\) in \(\varphi\) is free, and \emph{bound in} \(\varphi\) iff some
  occurrence of \(x\) in \(\varphi\) is bound.
\end{enumerate}
The prefixes of an address are linearly ordered by \(\sqsubseteq\), so a
bound occurrence is bound by exactly one quantifier occurrence: the longest
prefix of it that quantifies its variable.  Being free and being bound are
properties of occurrences; a variable can be both free and bound in one
formula.
\end{hbdefinition}

\begin{hbtheorem}{thm:p01-w02-p01-024-fv-occurrence}{proved=annex:p01-w02-p01-024-fv-occurrence}{The free-variable set collects the free occurrences}
For every first-order formula \(\varphi\),
\[
  \operatorname{Var}(\varphi)=\{\varphi[a]:a\in\operatorname{VO}(\varphi)\}
  \quad\text{and}\quad
  \FV(\varphi)=\{\varphi[a]:a\in\operatorname{VO}(\varphi)\text{ and }a
  \text{ is free in }\varphi\}.
\]
Consequently \(x\notin\FV(\varphi)\) iff every occurrence of \(x\) in
\(\varphi\) is bound, and \(\varphi\) is a sentence iff every variable
occurrence in it is bound.
\end{hbtheorem}

\begin{hbexample}{def:p01-w02-p01-024-binding-example}{One variable both free and bound}
In \(\varphi=P(x)\land\exists x\,Q(x,y)\) the variable occurrences are
\(1\cdot1\) (the \(x\) of \(P(x)\)), \(2\cdot0\) (the binding \(x\)),
\(2\cdot1\cdot1\) (the \(x\) of \(Q(x,y)\)) and \(2\cdot1\cdot2\) (the
\(y\)).  The only quantifier occurrence is \(2\); its scope is
\(\{2\cdot0,2\cdot1\cdot1,2\cdot1\cdot2\}\).  It binds \(2\cdot0\) and
\(2\cdot1\cdot1\), and does not bind \(2\cdot1\cdot2\) because it
quantifies \(x\), not \(y\).  So \(1\cdot1\) and \(2\cdot1\cdot2\) are free,
and \(x\) is both free and bound in \(\varphi\).  The recursion gives
\(\FV(\varphi)=\{x\}\cup(\{x,y\}\setminus\{x\})=\{x,y\}\), the variables of
the two free occurrences, as \cref{thm:p01-w02-p01-024-fv-occurrence}
requires.  In \(\forall x\,(P(x)\to\exists x\,Q(x))\) the occurrence
\(1\cdot2\cdot1\cdot1\) lies in the scope of both quantifier occurrences
\(\epsilon\) and \(1\cdot2\), and is bound by the inner one,
\(1\cdot2\).
\end{hbexample}

Precedence decides scope before binding is computed.  By the rank-3 prefix
rule of \cref{def:p01-w02-p01-026-precedence}, the abbreviation
\(\forall x\,P(x)\to Q(x)\) is the formula
\((\forall x\,P(x))\to Q(x)\), whose occurrence \(2\cdot1\) of \(x\) is
free; in \(\forall x\,(P(x)\to Q(x))\) the corresponding occurrence
\(1\cdot2\cdot1\) is bound.  Omitting the parentheses therefore changes the
free-variable set from \(\varnothing\) to \(\{x\}\).

\begin{hbwarning}{def:p01-w02-p01-024-binding-limits}{Scope of these definitions}
The definitions cover the binders \(\forall\) and \(\exists\) of the
first-order language.  A grammar with further binders (the description
operator \(\iota x\), the second-order quantifiers of
\cref{sec:prop-metatheory-01}, lambda abstraction, or the fixed-point
binders of the modal \(\mu\)-calculus) obtains its occurrence notions by
adding one quantifier-style clause per binder, with address \(0\) for the
bound variable; the theorem above is not asserted for such a grammar until
that grammar's free-variable recursion has been stated with the matching
clause.
\end{hbwarning}

\section{Association and fixity for the extended ranks}
\label{sec:p01-w02-p01-024-precedence}
\index{precedence!extended ranks}\index{association}

Table~\ref{tab:p01-precedence} lists every registry row that carries a
numerical rank, with its role and its recorded scope, and
\cref{def:p01-w02-p01-026-precedence} gives the fixity and the association
of the seven classical ranks and of the first nonclassical extensions.
The registry also ranks modal, temporal, epistemic, linear,
substructural, conditional and fuzzy operators beyond these.
Table~\ref{tab:p01-w02-p01-024-ranks} gives the fixity, the association
and the recorded scope of ranked rows outside the classical core; the
convention below says how operators that share a rank combine.  Thirty-one further registry rows of
Part~8 receive their ranks, or readings fixed by delimiters or by their
grammar, in \cref{def:p01-w02b-p01-012-p8-readings} and
Table~\ref{tab:p01-w02b-p01-012-p8-ranks}; twenty of them record the rank
given there and are listed in Table~\ref{tab:p01-precedence}, and eleven
carry no numerical rank.

\begin{hbconvention}{def:p01-w02-p01-024-fixity}{Fixity, association and the grammar boundary}
The convention applies only after the active language has been fixed and
its unique readability established, and only to operators whose recorded
scope contains the passage.  The grammar of a named primary calculus or
logic remains primary: where that grammar fixes a reading, the reading
stands, and where it requires parentheses, they are written.  In an
abbreviated expression, two operators \emph{compete} for an operand when the
operand is adjacent to both.
\begin{enumerate}
\item \emph{Rank.}  Of two competing operators, the one with the smaller
  rank takes the operand.
\item \emph{Same binary operator.}  If both competitors are occurrences of
  one binary operator, its association entry decides (in
  \cref{def:p01-w02-p01-026-precedence}, or in the association column of
  Table~\ref{tab:p01-w02-p01-024-ranks} or of
  Table~\ref{tab:p01-w02b-p01-012-p8-ranks}, and for the Part~9 program
  scale in Table~\ref{tab:p01-w02c-p01-017-ranks}): \emph{left} reads
  \(A\star B\star C\) as \((A\star B)\star C\), \emph{right} as
  \(A\star(B\star C)\), and \emph{none} licenses no reading, so
  parentheses must be written.
\item \emph{Different operators of equal rank.}  Two different operators of
  the same rank never compete in a licensed abbreviation: whichever fixities
  they have (two binary operators, a prefix and a postfix operator, or a
  prefix and a binary operator), parentheses must be written.  Hence
  \(!A^\perp\), \(A\otimes B\bindnasrepma C\), \(A\mathbin{\mathsf U}B
  \mathbin{\mathsf S}C\), \(\mathsf X^{I}A\mathbin{\mathsf U}B\) and
  \(A>B\to C\) are not abbreviations of anything.
\item \emph{Unary operators of one fixity.}  Two prefix operators, or two
  postfix operators, never compete, because nesting is forced:
  \(!?A\) is \(!(?A)\) and \(A^{\perp\perp}\) is \((A^\perp)^\perp\).
  A prefix operator otherwise takes the shortest following formula, and a
  postfix operator the shortest preceding one, subject to clauses 1 and 3.
\item \emph{Binders with a separating dot.}  For
  \(\lambda x.M\), \(\lambda x{:}\alpha.B\), \(\Lambda a.M\),
  \(\forall a.\alpha\), \(\mu Z.A\), \(\nu Z.A\), \(\exists r.C\) and
  \(\forall r.C\), the body after the dot is fixed by the source grammar.
  In the untyped and simply typed \(\lambda\)-calculi of
  \cref{sec:p07-lambda-calculi-and-type-theories} application binds more
  tightly than abstraction, so \(\lambda x.MN\) is \(\lambda x.(MN)\).
  Where the active grammar states no such rule, the binder is a rank-3
  prefix operator and clause 4 applies.
\item \emph{Bracketed and indexed forms.}  In \(O(A\mid B)\) the
  brackets and the bar delimit the two arguments completely: each argument
  is the whole expression between the delimiters, and the bracketed form is
  an indivisible operand outside.  An indexed conjunction or disjunction
  \(\bigwedge_{i\in I}\varphi_i\), \(\bigvee_{i\in I}\varphi_i\),
  \(\bigwedge\Gamma\) or \(\bigvee\Delta\) takes as operand only the
  displayed indexed formula or context, so
  \(\bigwedge_{i\in I}\varphi_i\land\psi\) is
  \((\bigwedge_{i\in I}\varphi_i)\land\psi\).
\item \emph{Term level and relations.}  Rank-1 operators act on terms and
  bind more tightly than every formula operator; postfix rank-1 signs such
  as \(N!\) and \(\kappa^+\) apply to the immediately preceding term.
  Rank-2 relations are nonassociative and combine terms into an atomic
  formula, so \(x<y<z\) is not an abbreviation of
  \(x<y\land y<z\).
\end{enumerate}
\end{hbconvention}

The program constructors of Part~9 use the separate program scale of
\cref{def:p01-w02c-p01-017-readings}: iteration is postfix at rank~1,
composition is left associative at rank~2, and choice is left associative
at rank~3.  These ranks are compared only with program ranks.  The same
convention gives the autoepistemic \(L\) a formula-prefix rank of~3;
prefix nesting fixes \(\neg LB=\neg(LB)\) and
\(L\neg B=L(\neg B)\).  Its scope and the delimiter rules for tests
and modalities are stated there.

\begingroup\footnotesize
\begin{longtable}{@{}p{0.05\textwidth}p{0.27\textwidth}p{0.11\textwidth}p{0.10\textwidth}p{0.31\textwidth}@{}}
\caption{Fixity and association of every registered rank beyond the
classical core.}
\label{tab:p01-w02-p01-024-ranks}\\
\toprule
Rank & Operators & Fixity & Association & Registry scope \\
\midrule
\endfirsthead
\caption[]{Fixity and association of every registered rank beyond the
classical core (continued).}\\
\toprule
Rank & Operators & Fixity & Association & Registry scope \\
\midrule
\endhead
\midrule\multicolumn{5}{r}{\emph{continued}}\\\endfoot
\bottomrule\endlastfoot
1 & \(N!\) & postfix & --- & P4 decision-hierarchy \\
1 & \(b^{\complement_{\mathbb B}}\) & postfix & --- & P6 cardinals-choice \\
1 & \(\kappa^+\) & postfix & --- & P6 cardinals-choice to large-cardinals-absoluteness \\
2 & \(\leq,<\) (canonical orderings) & infix & none & Parts 1--4 from set-prerequisites to universal-reductions; Legacy pre-batch5 P4.4-P4.6 reader meaning; Parts 4--19 and Parts 5--19 as recorded; P4 represented-space numerical ordering \\
2 & \(\leq,<\) (arithmetic order abbreviations) & infix & none & P4 decision-hierarchy; P5 arithmetic-coding \\
2 & \(>\) & infix & none & global \\
2 & \(\leq_{\N},<_{\N}\) & infix & none & Parts 4--5 from decision-hierarchy to arithmetic-coding; P4 decision-hierarchy \\
2 & \(<_{\mathcal O},\leq_{\mathcal O}\) & infix & none & Legacy pre-batch5 P4.5 reader meaning; P4 delta-one hyperarithmetical construction \\
2 & \(x\doteq y\) & infix & none & P9 legacy unit 4 temporal-automata \\
2 & \(q_1\subseteq q_2\) & infix & none & P9 legacy unit 7 logic-programming \\
2.5 & \(!A\) & prefix & --- & P8 paraconsistent-and-nonmonotonic \\
2.5 & \(?A\) & prefix & --- & P8 paraconsistent-and-nonmonotonic \\
2.5 & \(A^\perp\) & postfix & --- & P8 paraconsistent-and-nonmonotonic \\
2.5 & \(A^\circ\) & postfix & --- & P8 further-reading \\
3 & \(\Box,\Diamond\) & prefix & --- & Parts 5--8 from arithmetic-coding to conditionals-deontic-free \\
3 & \(\Box A\) (relevance abbreviation) & prefix & --- & P8 paraconsistent-and-nonmonotonic \\
3 & \(\mathsf X,\mathsf F,\mathsf G\) & prefix & --- & Parts 8--9 from conditionals-deontic-free \\
3 & \(K_i,E_G,C_G,B_i,D_G\) & prefix & --- & P8 conditionals-deontic-free \\
3 & \(\mathsf E,\mathsf A;\ \mathsf{EX},\mathsf{AX},\mathsf{EF},\mathsf{AF}\) & prefix & --- & Parts 8--9 from conditionals-deontic-free \\
3 & \(\mu Z.A,\nu Z.A\) & prefix binder & --- & Parts 8--9 from conditionals-deontic-free \\
3 & \(\langle a\rangle A,[a]A;\ [!\chi]A,\langle!\chi\rangle A;\ @_iA\) & prefix & --- & P8 conditionals-deontic-free \\
3 & \(U\,A\,B,\ S\,A\,B\) (Kamp) & prefix, two operands & --- & P8 conditionals-deontic-free \\
3 & \(\mathsf P,\mathsf H\) & prefix & --- & P8 conditionals-deontic-free \\
3 & \(A\alpha,\mathcal F\alpha\); \((\forall^E x),(\exists^E x)\) & prefix & --- & P8 modal-temporal-epistemic \\
3 & \(O A,\mathsf P_OA\) & prefix & --- & P8 many-valued-paraconsistent \\
3 & \(LA\) & prefix & --- & P8 further-reading \\
3 & \(\exists r.C,\forall r.C\) & prefix binder & --- & P9 legacy unit 3 description-logics \\
3 & \(\mathsf Y A\) & prefix & --- & P9 legacy unit 4 temporal-automata \\
3 & \(\mathsf P_{\bowtie q}[\psi]\) and its \(\min,\max\) forms; \(\mathsf S_{\bowtie q}[\Phi]\) & prefix, bracketed operand & --- & P9 legacy unit 5 model-checking \\
3 & \([\alpha]A\) & prefix & --- & P9 legacy unit 6 program-logics \\
3 & \(\mathsf{not}\ A\) & prefix & --- & P9 legacy unit 7 logic-programming \\
3 & \(E!(t)\); \(\iota x\,\varphi\) & predicate form; prefix binder & --- & Parts 2--8 from fol-syntax-semantics to many-valued-paraconsistent \\
3 & \(\mathsf{Eq}_{\alpha},\forall_{\alpha}\) & prefix & --- & P2 prop-metatheory \\
3 & \(\lambda x{:}\alpha.B\) & prefix binder & --- & Legacy pre-batch4 P2.3-P7.6 reader meaning \\
3 & \(\lambda x.M\); \(\forall a.\alpha\), \(\Lambda a.M\); \(\Pi_{x:A}B,\Sigma_{x:A}B\) & prefix binder & --- & P7 lambda-calculi-and-type-theories \\
3 & \(\alpha\mathbin{\times_{\mathrm{ty}}}\beta,\alpha\mathbin{+_{\mathrm{ty}}}\beta\) & infix & none & P7 normalization-curry-howard \\
3 & \(\varepsilon x\,\varphi(x)\) & prefix binder & --- & Legacy pre-batch4 P7.7 reader meaning \\
3.5 & \(A\mathbin{\mathsf U}B\) & infix & none & Parts 8--9 from conditionals-deontic-free \\
3.5 & \(A\mathbin{\mathsf S}B\) & infix & none & P9 legacy unit 4 temporal-automata \\
3.5 & \(\mathsf U^{\leq k},\mathsf U^{I}\); \(\mathsf X^{I}\) & infix; prefix & none; --- & P9 legacy unit 5 model-checking \\
4 & \(\otimes\), \(\bindnasrepma\), \(\mathbin{\&}\) (with) & infix & left & P8 paraconsistent-and-nonmonotonic \\
4 & \(A\mathbin{\cdot}B\) (Lambek product) & infix & left & P8 paraconsistent-and-nonmonotonic \\
4 & \(x\ast y\), \(\ast_{\text{\L}}\), \(\ast_{\Pi}\); \(A\mathbin{\&}B\) (strong conjunction) & infix & left & P8 relevance-linear \\
4 & \(P*Q\) (separating conjunction) & infix & left & P9 legacy unit 6 program-logics \\
4 & \(b\wedge_{\mathbb B}c\) & infix & left & P6 cardinals-choice \\
4 & \(\bigwedge_{i\in I}\varphi_i\) & indexed prefix & --- & P2 fol-syntax-semantics \\
4 & \(\bigwedge\Gamma\) & indexed prefix & --- & P7 lk-lj \\
5 & \(A\oplus B\) & infix & left & P8 paraconsistent-and-nonmonotonic \\
5 & \(b\vee_{\mathbb B}c\) & infix & left & P6 cardinals-choice \\
5 & \(\bigvee_{i\in I}\varphi_i\) & indexed prefix & --- & P2 fol-syntax-semantics \\
5 & \(\bigvee\Delta\) & indexed prefix & --- & P7 lk-lj \\
6 & \(\multimap\) & infix & right & P8 paraconsistent-and-nonmonotonic \\
6 & \(A\backslash B,\ B/A\) (Lambek residuals) & infix & none & P8 paraconsistent-and-nonmonotonic \\
6 & \(x\Rightarrow y\) (t-norm residuum) & infix & right & P8 relevance-linear \\
6 & \(A>B\) & infix & none & P8 many-valued-paraconsistent \\
6 & \(A\Box\!\!\to B\) & infix & none & P8 many-valued-paraconsistent \\
6 & \(O(A\mid B)\) & bracketed & --- & P8 many-valued-paraconsistent \\
6 & \(E(x)\supset\Phi\) & infix & right & P8 modal-temporal-epistemic \\
6 & \(\alpha\to_{\mathrm{ty}}\beta\) & infix & right & Legacy pre-batch4 P2.3-P7.6 reader meaning; P12 research-workflow \\
6 & \(C\sqsubseteq D\) & infix & none & P9 legacy unit 3 description-logics \\
\end{longtable}
\endgroup

\begin{hbexample}{def:p01-w02-p01-024-rank-examples}{Reading the extended ranks}
\begin{enumerate}
\item The \(\mathsf U\)-unfolding schema of the Hilbert calculus for linear
  temporal logic used in \cref{sec:p08-linear-temporal-completeness},
  \(A\mathbin{\mathsf U}B\leftrightarrow B\lor(A\land\mathsf X(A\mathbin{\mathsf U}B))\),
  reads as
  \((A\mathbin{\mathsf U}B)\leftrightarrow\bigl(B\lor(A\land\mathsf X(A\mathbin{\mathsf U}B))\bigr)\):
  until (3.5) binds more tightly than \(\lor\) (5) and \(\leftrightarrow\)
  (7), and the parentheses after \(\mathsf X\) are needed because the prefix
  \(\mathsf X\) (3) would otherwise take only \(A\).  Dropping the inner
  parentheses around \(A\land\mathsf X(\cdots)\) would not change the
  reading, because \(\land\) (4) binds more tightly than \(\lor\) (5).
\item \(\mathsf X A\mathbin{\mathsf U}B\lor C\) is
  \(((\mathsf X A)\mathbin{\mathsf U}B)\lor C\), whereas
  \(A\mathbin{\mathsf U}B\mathbin{\mathsf U}C\) is not licensed: until is
  nonassociative.
\item In linear logic \(A\bindnasrepma B\oplus C\) is
  \((A\bindnasrepma B)\oplus C\) and \(A\mathbin{\&}B\multimap C\) is
  \((A\mathbin{\&}B)\multimap C\); \(A\multimap B\multimap C\) is
  \(A\multimap(B\multimap C)\); \(?A\otimes B\) is \((?A)\otimes B\).
  The expressions \(!A^\perp\) and \(A\otimes B\bindnasrepma C\) must be
  written \(!(A^\perp)\) or \((!A)^\perp\), and
  \((A\otimes B)\bindnasrepma C\) or \(A\otimes(B\bindnasrepma C)\); the
  choice matters, since the two readings of each are different formulas.
\item In the da Costa scope, \(\neg A^\circ\) is \(\neg(A^\circ)\), because
  the postfix consistency sign (2.5) binds more tightly than negation (3).
\item In conditional logic \(A>(B>C)\) and \((A>B)>C\) must be written
  with their parentheses, and so must \(A>(B\to C)\).
\end{enumerate}
\end{hbexample}

\begin{hbwarning}{def:p01-w02-p01-024-rank-limits}{What the ranks do not decide}
A rank is a reading convention for abbreviations inside one declared
grammar.  It does not compare operators from disjoint scopes, it does not
make a formula of one logic a formula of another, and it gives no reading
to a parenthesis-free expression excluded by clauses 2 and 3 of
\cref{def:p01-w02-p01-024-fixity}.  Kamp's prefix connectives \(U\,A\,B\)
and \(S\,A\,B\) at rank 3 and the infix operators
\(\mathsf U,\mathsf S\) at rank 3.5 are different notations, recorded in
different scopes.
\end{hbwarning}

\section{The atomic diagram and the positive diagram}
\label{sec:p01-w02-p01-024-diagrams}
\index{diagram!atomic}\index{diagram!positive}

The model-theory lookup Table~\ref{tab:p03-model-symbols} records
\(\operatorname{Diag}(\mathcal M)\) with the meaning fixed by
\cref{def:diagrams}: the set of all atomic \emph{and negated atomic}
\(\Lang(M)\)-sentences true in \(\mathcal M_M\).  A smaller set is often
needed and must not be confused with it.

\begin{hbdefinition}{def:p01-w02-p01-024-positive-diagram}{Positive diagram}
For an \(\Lang\)-structure \(\mathcal M\) with domain \(M\), the
\emph{positive diagram} \(\operatorname{Diag}^+(\mathcal M)\) is the set of
all atomic \(\Lang(M)\)-sentences true in \(\mathcal M_M\), with the
parameter constants \(c_a\) of \cref{def:parameter-expansion}.  Thus
\[
  \operatorname{Diag}(\mathcal M)=\operatorname{Diag}^+(\mathcal M)\cup
  \{\neg\theta:\theta\text{ an atomic }\Lang(M)\text{-sentence with }
  \mathcal M_M\not\entails\theta\},
\]
and the second set contains \(c_a\neq c_b\) for all distinct
\(a,b\in M\), and \(\neg R(c_{a_1},\ldots,c_{a_n})\) whenever
\(R^{\mathcal M}\) omits \((a_1,\ldots,a_n)\).
\end{hbdefinition}

\begin{hbtheorem}{thm:p01-w02-p01-024-positive-diagram}{proved=annex:p01-w02-p01-024-positive-diagram}{What each diagram characterizes}
Let \(\mathcal M\) and \(\mathcal N\) be \(\Lang\)-structures.
\begin{enumerate}
\item There is a homomorphism \(h:\mathcal M\to\mathcal N\) iff
  \(\mathcal N\) has an \(\Lang(M)\)-expansion \(\mathcal N'\) with
  \(\mathcal N'\entails\operatorname{Diag}^+(\mathcal M)\).
\item There is an embedding \(h:\mathcal M\to\mathcal N\) iff
  \(\mathcal N\) has an \(\Lang(M)\)-expansion \(\mathcal N'\) with
  \(\mathcal N'\entails\operatorname{Diag}(\mathcal M)\)
  (the first part of \cref{thm:diagram-tests}).
\end{enumerate}
In both parts the map is \(h(a)=c_a^{\mathcal N'}\) and, conversely, the
expansion interprets \(c_a\) as \(h(a)\).
\end{hbtheorem}

\begin{hbexample}{def:p01-w02-p01-024-diagram-example}{A positive diagram without injectivity}
Let \(\Lang=\{<\}\), let \(\mathcal M=(\{0,1\},<)\) with \(0<1\), and let
\(\mathcal N\) be a one-element structure \((\{\ast\},<^{\mathcal N})\)
with \(\ast<^{\mathcal N}\ast\).  Interpreting both \(c_0\) and \(c_1\) as
\(\ast\) satisfies every atomic sentence: each equation \(c_a=c_b\) and each
instance of \(<\) is true at \(\ast\).  Hence
\(\mathcal N'\entails\operatorname{Diag}^+(\mathcal M)\), and the constant
map is a homomorphism.  But \(c_0\neq c_1\) and \(\neg(c_1<c_0)\) belong to
\(\operatorname{Diag}(\mathcal M)\) and fail in \(\mathcal N'\); no
expansion of \(\mathcal N\) satisfies \(\operatorname{Diag}(\mathcal M)\),
since a one-element domain cannot interpret \(c_0\) and \(c_1\) as distinct
elements.  Reading \(\operatorname{Diag}\) as the positive diagram would
therefore make the embedding test of \cref{thm:diagram-tests} false.
\end{hbexample}

\section{Single connectives, constants and Post's criterion}
\label{sec:p01-w02-p01-024-post}
\index{functional completeness!positive arity}\index{Post's criterion}
\index{clone!maximal}

\Cref{sec:prop-syntax-semantics-04} calls a connective set functionally
complete when it represents every \(f:\mathbb B^n\to\mathbb B\) for every
\(n\in\N\), and \(\N\) contains \(0\); a nullary function is represented
only by a formula with no variable.  Throughout this section truth values are written \(0,1\) as in
Table~\ref{tab:p01-sixteen}, and a formula over a connective set is built
from propositional variables by the connectives of the set only.  In this
section \(\operatorname{NAND}\) and \(\operatorname{NOR}\) are binary
connective symbols interpreted by rows 1110 and 1000 of that table, the
functions named in \cref{sec:p01-sixteen-functions}; they are connectives of
the formula sets studied here, not abbreviations within the classical
language.

\begin{hbconvention}{def:p01-w02-p01-024-arity-convention}{Positive-arity and nullary-inclusive completeness}
A set \(C\) of truth-functional connectives is \emph{positive-arity
functionally complete} iff for every \(n\geq1\) and every
\(f:\{0,1\}^n\to\{0,1\}\) some formula over \(C\) in the variables
\(p_1,\ldots,p_n\) has truth function \(f\).  It is
\emph{nullary-inclusive functionally complete} iff the same holds for every
\(n\in\N\); for \(n=0\) this requires a formula with no variable, that is,
a closed formula, for each of the two constants.  Functional completeness
in the sense of \cref{sec:prop-syntax-semantics-04} is the
nullary-inclusive notion.  The two nullary truth functions are represented
only through constants that the language declares, such as
\(\bottom\) and \(\topv\); a connective set with no nullary member gains no
nullary term formation from its positive-arity completeness.
\end{hbconvention}

\begin{hbtheorem}{thm:p01-w02-p01-024-nand-nor}{proved=annex:p01-w02-p01-024-nand-nor}{NAND and NOR are complete for positive arities}
Each of \(\{\operatorname{NAND}\}\) and \(\{\operatorname{NOR}\}\) is
positive-arity functionally complete.  Each of
\(\{\operatorname{NAND},\bottom\}\) and \(\{\operatorname{NOR},\bottom\}\)
is nullary-inclusive functionally complete.
\end{hbtheorem}

\begin{hbtheorem}{thm:p01-w02-p01-024-no-closed}{proved=annex:p01-w02-p01-024-no-closed}{No nullary representation without a constant}
\begin{enumerate}
\item If \(C\) contains no nullary connective, then every formula over \(C\)
  contains a propositional variable.  Hence \(C\) represents neither
  nullary truth function, and \(C\) is not functionally complete in the
  nullary-inclusive sense; in particular neither \(\{\operatorname{NAND}\}\)
  nor \(\{\operatorname{NOR}\}\) is.
\item \(C\) is nullary-inclusive functionally complete iff \(C\) is
  positive-arity functionally complete and contains at least one nullary
  connective.
\end{enumerate}
\end{hbtheorem}

So the sentence ``the Sheffer stroke and the Peirce arrow are each
functionally complete'' is true for positive arities and false in the
nullary-inclusive sense; nullary functions need the separately declared
constants or syntax.  The same convention governs Post's classification,
which, as Post states in his introduction, excludes functions of no
variables, that is, constants
\cite[Introduction, p.~2, n.~5]{Post1941}.

\begin{hbdefinition}{def:p01-w02-p01-024-clones}{Clones and the five Post classes}
For \(n\geq1\) let \(O_n\) be the set of functions
\(\{0,1\}^n\to\{0,1\}\) and \(O=\bigcup_{n\geq1}O_n\).  The \(i\)-th
\(n\)-ary projection is \(e^n_i(a_1,\ldots,a_n)=a_i\).  For
\(f\in O_m\) and \(g_1,\ldots,g_m\in O_n\), the \emph{composition}
\(f(g_1,\ldots,g_m)\in O_n\) sends \(\bar a\) to
\(f(g_1(\bar a),\ldots,g_m(\bar a))\).  A \emph{clone} is a subset of
\(O\) that contains every projection and is closed under composition; the
clone \(\langle C\rangle\) \emph{generated} by \(C\subseteq O\) is the
intersection of all clones containing \(C\).  A clone is \emph{proper} when
it is not \(O\), and \emph{maximal} when it is proper and the only clone
properly including it is \(O\).  For \(f\in O_n\) write
\(\bar0,\bar1\) for the constant tuples, \(\leq\) for the coordinatewise
order on \(\{0,1\}^n\), and \(\bar a^{\,c}\) for the coordinatewise
complement.  The five classes are:
\begin{enumerate}
\item \(\mathsf T_0=\{f:f(\bar0)=0\}\), the \(0\)-preserving functions;
\item \(\mathsf T_1=\{f:f(\bar1)=1\}\), the \(1\)-preserving functions;
\item \(\mathsf S=\{f:f(\bar a^{\,c})=1-f(\bar a)\text{ for all }\bar a\}\),
  the self-dual functions;
\item \(\mathsf M=\{f:\bar a\leq\bar b\text{ implies }f(\bar a)\leq
  f(\bar b)\}\), the monotone functions;
\item \(\mathsf L=\{f:f(\bar a)=c_0\oplus c_1a_1\oplus\cdots\oplus c_na_n
  \text{ for some }c_0,\ldots,c_n\in\{0,1\}\}\), the affine functions,
  where \(\oplus\) is addition modulo \(2\).
\end{enumerate}
\end{hbdefinition}

\begin{hbtheorem}{thm:p01-w02-p01-024-post}{proved=annex:p01-w02-p01-024-post}{Post's criterion for positive arities}
Let \(C\subseteq O\) be a set of positive-arity truth functions.
\begin{enumerate}
\item Each of \(\mathsf T_0,\mathsf T_1,\mathsf S,\mathsf M,\mathsf L\) is a
  proper clone, and none of the five is included in another.
\item \(\langle C\rangle=O\) iff, for each of the five classes, some member
  of \(C\) lies outside it.
\item Every proper clone is included in one of the five classes, and the
  maximal clones are exactly \(\mathsf T_0,\mathsf T_1,\mathsf S,\mathsf M\)
  and \(\mathsf L\).
\end{enumerate}
\end{hbtheorem}

The convention needed to apply the theorem to a language with constants is
explicit: \(\langle C\rangle\) consists of positive-arity functions only, so
a nullary constant symbol \(\bottom\) or \(\topv\) of a language enters the
criterion as the unary constant function \(p\mapsto0\) or \(p\mapsto1\), and
the conclusion \(\langle C\rangle=O\) is positive-arity completeness of the
language's connectives (the truth functions of formulas over a connective
set are exactly the members of the clone its functions generate, by the
correspondence used in the annex).  Nullary-inclusive completeness then
holds exactly when the language also declares a nullary constant, by the
second part of \cref{thm:p01-w02-p01-024-no-closed}.

Post proves the classification for iteratively closed systems of two-valued
functions.  In his notation the complete system is \(C_1\); the five
systems immediately included in it are \(C_2\), \(C_3\), \(D_3\), \(A_1\)
and \(L_1\); every closed system properly included in \(C_1\) is included
in at least one of them; and none of the five is properly included in
another \cite[\S\,25, p.~99]{Post1941}.  A set of generators therefore
generates the complete system exactly when each of the five fails to
contain it \cite[\S\,26, p.~105]{Post1941}.  \(D_3\) is the system of all
self-dual functions, \(A_1\) the system of functions satisfying his
\([A{:}a]\) condition and \(L_1\) the system of all alternating functions
\cite[\S\,26, p.~106]{Post1941}.  Post classifies a function by the
first-order function obtained by identifying all its arguments: an
\(\alpha\)-, \(\beta\)-, \(\gamma\)- or \(\delta\)-function according as
that function is the identity, the constant \(1\), the constant \(0\) or
the negation \cite[\S\,10, p.~43]{Post1941}.  His statement that a set
escapes \(C_2\) only through a \(\gamma\)- or \(\delta\)-function, and
escapes \(C_3\) only through a \(\beta\)- or \(\delta\)-function
\cite[\S\,26, p.~105]{Post1941}, says that \(C_2\) and \(C_3\) play the
roles of \(\mathsf T_1\) and \(\mathsf T_0\).  Since every
\(\delta\)-function fails the \([A{:}a]\) condition
\cite[\S\,26, p.~106]{Post1941} while the negation, a \(\delta\)-function,
lies in \(\mathsf L\), the class \(A_1\) is not \(\mathsf L\); as both lists
consist of exactly the maximal clones, \(A_1\) is \(\mathsf M\) and \(L_1\)
is \(\mathsf L\).  The proof in the annex is self-contained and does not
depend on this comparison.

\begin{hbexample}{def:p01-w02-p01-024-post-example}{Applying the criterion}
\begin{enumerate}
\item NAND is outside all five classes: \(\operatorname{NAND}(0,0)=1\),
  \(\operatorname{NAND}(1,1)=0\); \(\operatorname{NAND}(0,1)=1\) while
  \(\operatorname{NAND}(1,0)=1\), so it is not self-dual; it decreases from
  \((0,0)\) to \((1,1)\); and \(\operatorname{NAND}(a,b)=1\oplus ab\) is not
  affine, because the affine functions of two variables take the value \(1\)
  at an even number (zero, two or four) of the four inputs, and NAND takes
  it at three.  This gives \cref{thm:p01-w02-p01-024-nand-nor} again for
  positive arities; NOR is handled by the same five checks.
\item \(\{\to\}\) is included in \(\mathsf T_1\), since \(1\to1=1\), so it
  is not complete.  With the constant \(\bottom\), entered as the unary
  function \(p\mapsto0\), the set escapes \(\mathsf T_1\); \(\to\) escapes
  \(\mathsf T_0\) (\(0\to0=1\)), \(\mathsf M\) (\((0,0)\leq(1,0)\) but
  \(1>0\)), \(\mathsf S\) (\(0\to0=1=1\to1\)) and \(\mathsf L\) (it takes
  the value \(1\) at three inputs).  So \(\{\to,\bottom\}\) is positive-arity
  complete, and nullary-inclusive complete as well, because
  \(\bottom\to\bottom\) is a closed formula for the constant \(1\).
\item \(\{\land,\lor,\bottom,\topv\}\), with the constants entered as unary
  constant functions, lies inside \(\mathsf M\), so it does not define
  negation; \(\{\neg,\leftrightarrow\}\) lies inside \(\mathsf L\), so it
  does not define conjunction.
\end{enumerate}
\end{hbexample}

\begin{hbwarning}{def:p01-w02-p01-024-post-limits}{Limits of the criterion}
The criterion concerns composition of functions, not derivability: a
complete connective set need not come with a complete calculus.  It is
specific to two values; Post's classification of the closed systems is
for two-valued functions only, and the five-class test does not transfer to
many-valued truth functions.  It also concerns positive arities only, as
the convention above requires.
\end{hbwarning}

\section{Relationships}
\label{sec:p01-w02-p01-024-relationships}

The letter-family tables refine Table~\ref{tab:p01-letter-families} with the
complete Fraktur and blackboard rows and with every Greek metavariable
family; the master register of \cref{sec:p01-generated-symbols} remains the
row-level authority.  The occurrence definitions supply the binding notion
used by substitution and \(\alpha\)-conversion in
\cref{sec:fol-syntax-semantics-03} and later parts.  The fixity convention,
with \cref{def:p01-w02-p01-026-precedence}, completes
Table~\ref{tab:p01-precedence}.  The positive diagram is the
homomorphism counterpart of the embedding and elementary-embedding tests of
\cref{thm:diagram-tests}.  The completeness results refine the basis theorem
of \cref{sec:prop-syntax-semantics-04} and the definability notes after
Table~\ref{tab:p01-sixteen}.
